\documentclass[10pt,a4paper]{amsart}
\usepackage[margin=19mm]{geometry}
\usepackage{amsmath,amssymb,mathtools}
\usepackage[scr=rsfs,cal=euler]{mathalfa}
\usepackage{xfrac}
\usepackage[unicode,hidelinks,bookmarks=false]{hyperref}
\newcommand{\ie}{i.e.\ }
\newcommand{\cf}{cf.\ }
\newcommand{\resp}{respectively\ }
\newcommand{\Subsection}{\subsection}
\providecommand{\nobreakdash}{\nobreak\hbox{-}}
\setlength{\emergencystretch}{2em}
\multlinegap0pt
\usepackage[shortlabels]{enumitem}
\newtheorem{theorem}{Theorem}
\newtheorem{proposition}{Proposition}
\newtheorem{lemma}{Lemma}
\title{Symmetry of bounded solutions to quasilinear elliptic equations in a half-space}
\author{Phuong Le}
\date{Source: arXiv:2409.04804v1 (CC BY 4.0)}
\begin{document}
\maketitle

\noindent\emph{Source and licence.} Symmetry of bounded solutions to quasilinear elliptic equations in a half-space. Version arXiv:2409.04804v1 (CC BY 4.0), distributed by arXiv under the Creative Commons Attribution 4.0 International licence, http://creativecommons.org/licenses/by/4.0/, as recorded in the arXiv licence metadata for this exact version and shown on the abstract page https://arxiv.org/abs/2409.04804v1. Full licence notice: \texttt{licenses/arxiv-2409.04804-CC-BY-4.0.txt}.\par\smallskip

\noindent\textit{Editorial scope.} Counted unit: the article's lemma lem:lower (source0.tex lines 658--660). The article's complete original proof of it is delivered with it (source0.tex lines 661--815, one \texttt{proof} environment of 155 lines). No mathematics is written, repaired or re-derived: the statement and the proof are the article's own lines, in the article's own notation, and no retained line is substituted or re-flowed.  Retained uncounted as the source-local context the proof needs: lines 57--63; lines 109--126; lines 111--115; lines 143--149; lines 152--180; lines 239--247; lines 250--260; lines 263--273; lines 276--295; lines 486--492; lines 494--504; lines 617--619.  Nothing was omitted from any retained span, so the package carries no omission record.  No retained line contains a citation, so the wrapper carries no bibliography and no bibliography entry.  No retained line contains a figure environment, an image-inclusion command or drawing code, so no image is delivered and the package declares no asset dependency.  Editorial formatting only, in addition: an AMS \texttt{amsart} preamble that restates the article's own declarations and macros used by the retained lines, namely the environments \texttt{theorem}, \texttt{proposition}, \texttt{lemma} on separate counters, together with the standard packages those lines need.  The retained environments print in this document as Theorem 1, Proposition 1, Lemma 1 in order of appearance, whereas the article prints the same items with the numbering of its own full document.  The complete native source of the article as distributed in the arXiv e-print for this exact version is bundled unchanged as \texttt{sources/164919/source0.tex}.
	\begin{equation}\label{main}
		\begin{cases}
			-\Delta_p u = f(u) &\text{ in } \mathbb{R}^N_+,\\
			u > 0 &\text{ in } \mathbb{R}^N_+,\\
			u = 0 &\text{ on } \partial\mathbb{R}^N_+,
		\end{cases}
	\end{equation}

	\begin{theorem}\label{th:main}
		Assume $p>1$ and $f:[0,+\infty)\to\mathbb{R}$ is a continuous function which is locally Lipschitz continuous in $[0,+\infty)\setminus Z_f$. Let $u\in C^1(\overline{\mathbb{R}^N_+})$ be a bounded solution to \eqref{main}. Assume that $\rho=\sup_{\mathbb{R}^N_+} u$ verifies $f(\rho)=0$ and $f$ has no zeros in $(\rho-\varepsilon,\rho)$ for some $\varepsilon>0$. Moreover, assume that		
		\begin{equation}\label{smp_cond}
			\liminf_{t\to z^+} \frac{f(t)}{(t-z)^{p-1}} > -\infty,\quad
			\limsup_{t\to z^-} \frac{f(t)}{(z-t)^{p-1}} < +\infty
			\quad\text{ for all } z\in Z_f
		\end{equation}
		and one of the following conditions hold
		\begin{enumerate}
			\item[(i)] $f(0)\ge0$,
			\item[(ii)] $F(\rho)\ne0$,
			\item[(iii)] $u\in C^2(\mathbb{R}^N_+\cap B_r(x_0))$ for some $x_0\in\partial\mathbb{R}^N_+$ and $r>0$.
		\end{enumerate}
		Then $u$ depends only on $x_N$ and monotone increasing in $x_N$. Furthermore,
		\begin{equation}\label{deri}
			\frac{\partial u}{\partial x_N} > 0 \text{ in } \mathbb{R}^N_+.
		\end{equation}
	\end{theorem}

		\begin{equation}
			\liminf_{t\to z^+} \frac{f(t)}{(t-z)^{p-1}} > -\infty,\quad
			\limsup_{t\to z^-} \frac{f(t)}{(z-t)^{p-1}} < +\infty
			\quad\text{ for all } z\in Z_f
		\end{equation}

	\begin{equation}\label{1D}
		\begin{cases}
			-\Delta_p u = f(u) & \text{ in } \mathbb{R}_+,\\
			u \ge 0 & \text{ in } \mathbb{R}_+,\\
			u(0)=0, ~ \|u\|_{L^\infty(\mathbb{R}_+)}<+\infty,
		\end{cases}
	\end{equation}

	\begin{theorem}\label{th:1D}
		Assume $p>1$ and $f:[0,+\infty)\to\mathbb{R}$ is a continuous function which is locally Lipschitz continuous in $[0,+\infty)\setminus Z_f$ such that \eqref{smp_cond} holds. Then we have the following claims:
		\begin{enumerate}
			\item[(i)] For each $\rho\in Z_f^*:=\{z\in Z_f \mid F(t) < F(z) \text{ for all } t\in [0,z)\} \cup Z_f^0$, where			
			\[
			Z_f^0=
			\begin{cases}
				\{z\in Z_f \mid F(t) < 0 = F(z) \text{ for all } t\in (0,z)\} &\text{ if } f(0) < 0,\\
				\emptyset &\text{ otherwise},
			\end{cases}
			\]		
			there exists a unique solution $u_\rho$ to \eqref{1D} such that $u_\rho'>0$ in $\mathbb{R}_+$, $\lim\limits_{t\to+\infty} u_\rho(t)=\rho$ and $u_\rho'(0)=\left(\frac{p}{p-1}F(\rho)\right)^\frac{1}{p}$.
			\item[(ii)] If $f(0)<0$, then for
			\[
			\rho\in P_f := \{z>0 \mid F(t) < 0 = F(z) \text{ for all } t\in (0,z) \text{ and } f(z)>0\},
			\]
			there exists a unique solution $\tilde u_\rho$ to \eqref{1D} such that $\tilde u_\rho'>0$ in $(0,t^*)$, $\tilde u_\rho'(0)=\tilde u_\rho'(t^*)=0$, $\tilde u_\rho(t^*)=\rho$ and
			\[
			\tilde u_\rho(t+2kt^*)=\tilde u_\rho(t), \quad \tilde u_\rho(t+(2k+1)t^*)=\tilde u_\rho(t^*-t)
			\]
			for all $t\in[0,t^*]$ and all $k\in\mathbb{N}$,
			where
			\[
			t^* = \left(\frac{p-1}{p}\right)^\frac{1}{p}\int_{0}^{\rho} \frac{ds}{\left[F(\rho)-F(s)\right]^\frac{1}{p}}.
			\]
		\end{enumerate}
		
		Moreover, every solution $u \not\equiv 0$ to \eqref{1D} must have one of the above two forms.
	\end{theorem}

	\begin{theorem}[Strong maximum principle and Hopf's lemma]\label{th:smp}
		Let $u\in C^1(\Omega)$ be a nonnegative weak solution to
		\[
		-\Delta_p u + cu^q = g \ge 0 \quad\text{ in }\Omega,
		\]
		where $p>1$, $q\ge p-1$, $c\ge0$, $\Omega$ is a connected domain of $\mathbb{R}^N$ and $g \in L^\infty_{\rm loc}(\Omega)$. If $u \not\equiv 0$ then $u > 0$ in $\Omega$. Moreover for any point $x_0 \in \partial \Omega$ where the interior sphere condition is satisfied, and such that $u\in C^1(\Omega\cup\{x_0\})$ and $u(x_0) = 0$ we have that $\frac{\partial u}{\partial \nu}>0$ for any inward directional
		derivative (this means that if $x$ approaches $x_0$ in a ball $B\subset\Omega$
		that has $x_0$ on its boundary, then $\lim\limits_{x\to x_0} \frac{u(x)-u(x_0)}{|x-x_0|}>0$).
	\end{theorem}

	\begin{theorem}[Strong comparison principle]\label{th:scp}
		Let $u,v\in C^1(\Omega)$ be two solutions to
		\[
		-\Delta_p w = f(w) \quad\text{ in } \Omega
		\]
		such that $u \le v$ in $\Omega$, with $p>1$ and let
		\[
		Z=\{x\in\Omega \mid |\nabla u(x)|+|\nabla v(x)|=0\}.
		\]
		If $x_0\in\Omega\setminus Z$ and $u(x_0)=v(x_0)$, then $u=v$ in the connected component of $\Omega\setminus Z$ containing $x_0$.
	\end{theorem}

	\begin{theorem}[Boundary point lemma]\label{lem:bpm}
		Let $u,v\in C^1(\overline{\Omega})\cap C^2(\Omega)$ be solutions of the equation
		\[
		-\Delta_p w = f(w) \quad\text{ in }\Omega
		\]
		such that $u<v$ in $\Omega$ and $u=v$ at some point $x_0 \in \partial \Omega$ where the interior sphere condition is satisfied. Then
		\[
		\frac{\partial u}{\partial \nu}(x_0)<\frac{\partial v}{\partial \nu}(x_0),
		\]
		where $\nu$ is the inward normal at $x_0$.
	\end{theorem}

	\begin{theorem}[Weak sweeping principle]\label{th:wsp}
		Suppose that $\Omega$ is a bounded smooth domain in $\mathbb{R}^N$, $h(x, s)$ is measurable in $x \in \Omega$,
		continuous in $s$, and for each finite interval $J$, there exists a continuous increasing function $L(s)$ such
		that $h(x, s) + L(s)$ is nondecreasing in $s$ for $s \in J$ and $x \in \Omega$. Let $u_t$ and $v_t$, $t \in [t_1, t_2]$, be functions in $W^{1,p}(\Omega) \cap C(\overline \Omega)$ and satisfy in the weak sense,
		\[
		\begin{cases}
			-\Delta_p u_t \ge h(x,u_t) + \varepsilon_1(t) & \text{ in } \Omega,\\
			-\Delta_p v_t \le h(x,v_t) - \varepsilon_2(t) & \text{ in } \Omega,\\
			u_t \ge v_t + \varepsilon & \text{ on } \partial\Omega,\\
		\end{cases}
		\]
		for all $t \in [t_1, t_2]$, where
		\[
		\varepsilon_1(t) + \varepsilon_2(t) \ge \varepsilon > 0.
		\]
		Moreover, suppose that $u_{t_0} \ge v_{t_0}$ in $\Omega$ for some $t_0 \in [t_1, t_2]$ and $t \mapsto u_t$, $t \mapsto v_t$ are continuous from the finite closed interval $[t_1, t_2]$ to $C(\overline \Omega)$. Then
		\[
		u_t \ge v_t \text{ in } \Omega \text{ for all } t \in [t_1, t_2].
		\]
	\end{theorem}

	\begin{equation}\label{ball}
		\begin{cases}
			-\Delta_p u = f(u) &\text{ in } B_r,\\
			u > 0 &\text{ in } B_r,\\
			u = 0 &\text{ on } \partial B_r,
		\end{cases}
	\end{equation}

	\begin{proposition}\label{prop:ball}
		Assume that $f$ is a locally Lipschitz continuous function with $f (0) \ge 0$ and $f(\rho) = 0$ for some $\rho > 0$. Assume in addition that
		\[
		F_\rho(t):=\int_{t}^{\rho} f(s) ds > 0 \quad\text{ for all } t \in [0, \rho).
		\]
		Then for every $\varepsilon > 0$ there exists a positive number
		$R_0 = R_0(\varepsilon)$ such that for $r \ge R_0$, problem \eqref{ball} admits a solution $u_r \in C^{1,\alpha}(\overline{B_r})$ such that
		\[
		\rho - \varepsilon \le \|u_r\|_{L^\infty(B_r)} < \rho.
		\]
	\end{proposition}

	\begin{lemma}\label{lem:upper}
		Under assumptions of Theorem \ref{th:main}, there exists a one-dimensional solution $\overline u$ of \eqref{main} such that $u \le \overline u \le \rho$.
	\end{lemma}

	\begin{lemma}\label{lem:lower}
		Under assumptions of Theorem \ref{th:main}, there exists a one-dimensional solution $\underline u$ of \eqref{main} such that $\|\underline u\|_{L^\infty(\mathbb{R}^N)}=\rho$ and $\underline u \le u$ in $\mathbb{R}^N$.
	\end{lemma}

	\begin{proof}
		For $\alpha<\beta$, we denote
		\[
		\Sigma_{\alpha,\beta} = \{x\in\mathbb{R}^N \mid \alpha < x_N < \beta\} \quad\text{ and }\quad \Sigma_\alpha = \Sigma_{\alpha,\infty} := \{x\in\mathbb{R}^N \mid \alpha < x_N \}.
		\]
		
		We will construct lower solution $\underline u$ in 3 steps:
		
		\textit{Step 1.} For every $r > 0$ and every $\varepsilon > 0$ small enough, there exists $x_0 \in \mathbb{R}^N_+$
		such that $B_r(x_0) \subset \overline\Sigma_1$ and $u \ge \rho - \varepsilon$ in $B_r(x_0)$.
		
		To prove this assertion we take a sequence of points $x_n=(x_n', x_{n,N}) \in \mathbb{R}^N_+$ such that $\lim_{n\to\infty} u(x_n) = \rho$. Since the right hand side of \eqref{main} is bounded, by standard elliptic estimates, we have $u\in C^{1,\alpha}(\overline{\mathbb{R}^N_+})$. Because of the Dirichlet condition, this implies that $(x_{n,N})$ is bounded away from zero. Hence extracting a subsequence we may assume that either $x_{n,N}\to y_0$ for some $y_0>0$ or $x_{n,N}\to+\infty$. We define
		\[
		u_n(x) = u(x+x_n).
		\]
		
		Passing to a subsequence, we have $u_n\to v$ in $C^{1,\alpha'}(\overline{\mathbb{R}^N_+})$, where $v$ is a solution of
		\[
		\begin{cases}
			-\Delta_p v = f(v) & \text{ in } \Sigma,\\
			0 \le v \le \rho & \text{ in } \Sigma,
		\end{cases}
		\]
		where $\Sigma=\{x_N \ge -y_0\}$ in case $x_{n,N}\to y_0$ or $\Sigma=\mathbb{R}^N$ when $x_{n,N}\to+\infty$. In either case, using \eqref{smp_cond} and the fact that $f (\rho) = 0$ and $v(0)=\rho$, the strong maximum principle (Theorem \ref{th:smp}) implies $v\equiv\rho$. However, if $x_{n,N}\to y_0$, then from the Dirichlet condition on $u$, we have $v=0$ on $\partial \Sigma$, a contradiction.
		
		Therefore, the case $x_{n,N}\to+\infty$ must happen and $u(x+x_n) \to \rho$ locally uniformly in $\mathbb{R}^N$.
		
		Consequently, for any $\varepsilon > 0$ and $r > 0$, we have $u(x + x_n) \ge \rho - \varepsilon$ in $B_r$ if $n$ is larger than some $n_0 = n_0(\varepsilon, r)$. Then $u(x) \ge \rho - \varepsilon$ in $B_r(x_n)$ for those values of $n$.
		
		\textit{Step 2.} For every $\varepsilon > 0$, there exists $\eta > 0$ such that $u(x) \ge \rho - \varepsilon$ for $x\in\overline\Sigma_\eta$.
		
		By Lemma \ref{lem:upper}, problem \eqref{main} has a one-dimensional solution $\overline u$ such that $\sup_{\mathbb{R}^N_+} \overline u = \rho$ and $u \le \overline u$. Since $\frac{\partial\overline u}{\partial x_N}>0$, the strong comparison principle (Theorem \ref{th:scp}) can be applied to yield that either $u \equiv \overline u$ or $u < \overline u$. If the former case happens, then the lemma is proved by setting $\underline u = \overline u$.
		
		Hence in what follows, we assume $u < \overline u$ in $\mathbb{R}^N_+$. We claim that
		\begin{equation}\label{s1}
			F(t) < F(\rho) \quad\text{ for all } t\in[0,\rho).
		\end{equation}	
		Indeed, if (i) or (ii) holds, then \eqref{s1} follows from Theorem \ref{th:1D}.
		If (iii) holds, then both $u$ and $\overline u$ belong to $C^2(\mathbb{R}^N_+\cap B_r(x_0))$. Hence we can apply the boundary point lemma (Theorem \ref{lem:bpm}) to obtain
		\[
		\frac{\partial \overline u}{\partial x_N}>\frac{\partial u}{\partial x_N}\ge0.
		\]
		Therefore, an application of Theorem \ref{th:1D} also yields \eqref{s1}.
		
		If $f (0) < 0$, we choose $\mu>0$ small such that
		\begin{equation}\label{s2}
			\frac{f(0)}{2}\mu + F(\rho) > 0.
		\end{equation}
		Then we define function $f_\mu:[-\mu,\rho]\to\mathbb{R}$ by setting
		\[
		f_\mu(t) =
		\begin{cases}
			\frac{f(0)}{\mu}(t+\mu) & \text{ if } t\in[-\mu,0),\\
			f(t) & \text{ if } t\in[0,\rho].
		\end{cases}
		\]
		In the case $f (0) \ge 0$, we simply take $\mu = 0$ and $f_0 = f$. From \eqref{s1} and \eqref{s2}, we deduce
		\begin{equation}\label{s3}
			\int_{t}^{\rho} f_\mu(s) ds > 0 \quad\text{ for all } t\in[-\mu,\rho).
		\end{equation}
		
		Now for each small $\delta>0$, we define
		\[
		g_{\mu,\delta}(t) := f_\mu(t) - \delta (t+\mu) \quad\text{ in } [-\mu,\rho].
		\]
		
		Since $f$ has no zeros in a left neighborhood of $\rho$ and \eqref{s1} holds, we deduce that $f$ is positive in that neighborhood. On the other hand, $\int_{t}^{\rho} g_{\mu,\delta}(s) ds \to \int_{t}^{\rho} f_\mu(s) ds$ as $\delta\to0$, uniformly in $t\in[-\mu,\rho]$. Hence taking into account \eqref{s3}, for sufficiently small $\delta$, there exists $\rho_\delta\nearrow\rho$ (as $\delta\to0$) such that $g_{\mu,\delta}(\rho_\delta)=0$ and
		\[
		G(t):=\int_{t}^{\rho_\delta} g_{\mu,\delta}(s) ds > 0 \quad\text{ for all } t \in [-\mu, \rho_\delta).
		\]
		
		We can therefore apply Proposition \ref{prop:ball} to deduce that the problem
		\[
		\begin{cases}
			-\Delta_p w = g_{\mu,\delta}(w-\mu) &\text{ in } B_r,\\
			w = 0 &\text{ on } \partial B_r
		\end{cases}
		\]
		with some $r=r(\delta)$ has a positive solution $w_r$ verifying
		\[
		\rho_\delta + \mu - \delta \le \|w_r\|_{L^\infty(B_r)} < \rho_\delta + \mu.
		\]
		
		Hence by setting $v_r=w_r-\rho$, we have that $v_r\ge-\mu$,
		\[
		\rho_\delta - \delta \le \|v_r\|_{L^\infty(B_r)} < \rho_\delta
		\]
		and $v_r$ solves the problem
		\[
		\begin{cases}
			-\Delta_p v = g_{\mu,\delta}(v) &\text{ in } B_r,\\
			v = -\mu &\text{ on } \partial B_r.
		\end{cases}
		\]
		
		By Step 1, we can find $x_0$ so that
		$B_r(x_0) \subset \overline\Sigma_1$ and $u \ge \rho_\delta$ in $B_r(x_0)$.
		
		Let $\tilde{x}$ be an arbitrary point in $\overline\Sigma_{r+1}$. We define $x^t = (1-t)x^0 + t\tilde{x}$ and $u^t(x)=u(x+x^t)$.
		Clearly, $B_r(x^t) \subset \overline\Sigma_1$ for all $t \in [0, 1]$.
		
		Since $u > 0$ on the compact set $\bigcup_{t\in[0,1]} \overline{B_r(x^t)}$, we can find $\lambda > 0$ such that $u \ge \lambda$ on this set. Let $r_1 < r$ be such that $-\mu < v_r < \frac{\lambda}{2}$ on $\{|x| = r_1\}$. Then for all $t \in [0, 1]$, we have
		\[
		\begin{cases}
			-\Delta_p u^t = f(u^t) & \text{ in } B_{r_1},\\
			-\Delta_p v_r = g_{\mu,\delta}(v_r) \le f(v_r) - \gamma & \text{ in } B_{r_1},\\
			u^t \ge v_r + \frac{\lambda}{2} & \text{ on } \partial B_{r_1},
		\end{cases}
		\]
		where $\gamma=\inf_{x\in B_{r_1}} [f (v_r(x)) - g_{\mu,\delta}(v_r(x))] > 0$. Moreover, $u^0 = u \ge \rho_\delta \ge v_r$ in $B_{r_1}$. Thus we can apply the weak sweeping principle (Theorem \ref{th:wsp}) to deduce that $u^t \ge v_r$ in $B_{r_1}$ for all $t \in [0, 1]$. In particular, $u(\tilde{x}) \ge v_r(0) \ge \rho_\delta - \delta$. Since $\tilde{x} \in \overline\Sigma_{r+1}$ is arbitrary, we have
		\[
		u \ge \rho_\delta - \delta \quad \text{ in } \overline\Sigma_{r+1}.
		\]
		Since $\rho_\delta\to\rho$ as $\delta\to0$, we can choose $\delta$ small such that $\rho_\delta - \delta > \rho - \varepsilon$ for given $\varepsilon>0$. Hence the conclusion follows.
		
		\textit{Step 3.} There exists a minimal solution $\underline u$ of \eqref{main} such that $\|\underline u\|_{L^\infty(\mathbb{R}^N)}=\rho$ and $\underline u \le u$ in $\mathbb{R}^N$.
		
		Since $f$ has no zeros in a left neighborhood of $\rho$ and \eqref{s1} holds, there exists a small $\delta > 0$ such that $f > 0$ in $[\rho - \delta, \rho)$. Now we choose $\varepsilon<\delta$ sufficiently small and then modify $f$ over $[0, \rho]$ to obtain a new function $f_\varepsilon$ such that
		\[
		f_\varepsilon = f \text{ in } [0, \rho-\delta],\quad
		0<f_\varepsilon<f \text{ in } (\rho-\delta,\rho-\varepsilon),\quad
		f_\varepsilon(\rho-\varepsilon)=0
		\]
		and
		\[
		\int_{t}^{\rho-\varepsilon} f_\varepsilon(s) ds > 0 \quad\text{ for all } t\in[0, \rho-\varepsilon).
		\]
		
		Hence we can apply Theorem \ref{th:1D} to see that the problem
		\[
		\begin{cases}
			-\Delta_p v = f_\varepsilon(v) & \text{ in } \mathbb{R}_+,\\
			v(0)=0,~ \lim_{t\to+\infty} v(t) = \rho-\varepsilon
		\end{cases}
		\]
		has a unique positive solution $v_\varepsilon$, which is monotone increasing. For any $\eta>0$, we define
		\[
		v_{\varepsilon,\eta}(t) =
		\begin{cases}
			0 & t \in [0,\eta),\\
			v_\varepsilon(t-\eta) & t \in [\eta,+\infty).
		\end{cases}
		\]
		
		It is easy to check that $v_{\varepsilon,\eta}$ is a lower solution of \eqref{main}. By Step 2, we can find a large number $\eta$ so that $u>\rho-\varepsilon$ in $\{x_N\ge \eta\}$. Hence in view of $v_\varepsilon\le \rho-\varepsilon$, we find that $v_{\varepsilon,\eta}\le u$ in $\mathbb{R}^N_+$ for such $\eta$. For $n$ large, using the auxiliary problem
		\[
		\begin{cases}
			-\Delta_p u = f(u) &\text{ in } B_n^+,\\
			u = (\rho-\varepsilon)\chi_{\{x_N\ge\eta\}} &\text{ on } \partial B_n^+
		\end{cases}
		\]
		and arguing as in the proof of Lemma \ref{lem:upper}, we find that equation \eqref{main} has a minimal solution $\underline u$ in the order interval $[v_{\varepsilon,\eta}, u]$.
		
		Moreover, the minimality of $\underline u$ implies that $\underline u$ is a function of $x_N$ only. Thus $\underline u$ is a solution to \eqref{1D}. Due to the lower bound $v_{\varepsilon,\eta}$, we see that $\underline u$ is not periodic. Hence it must be positive and monotone by Theorem \ref{th:1D}. Moreover, $\lim_{t\to+\infty} \underline u(t)$ is a zero of $f$ contained in $[\rho-\varepsilon,\rho]$. But $\rho$ is the only zero of $f$ in this range. Hence $\lim_{t\to+\infty} \underline u(t)=\rho$. This implies $\|\underline u\|_{L^\infty(\mathbb{R}^N)}=\rho$.
	\end{proof}

\end{document}
